\section{Trazabilidad entre el levantamiento y el diseño}
\label{anexo:trazabilidad}

Esta tabla es la que permite comprobar, sin leer todo el documento, que cada requisito nace de algo
que el entrevistado dijo y que cada requisito tiene un lugar en la arquitectura.

\begin{table}[H]
\centering
\caption{Trazabilidad: pregunta del guion $\rightarrow$ requisito $\rightarrow$ contenedor}
\label{tab:trazabilidad}
\small
\begin{tabularx}{\linewidth}{@{}llYl@{}}
\toprule
Pregunta del guion & Requisito & Decisión de arquitectura que lo atiende & Contenedor \\
\midrule
Guion, P5 (sin validar) & RF-01 & Alta rápida y persistencia desacoplada del adjunto (ADR-002) & \texttt{SVC-INC} \\
\addlinespace
Guion, P6 (sin validar) & RF-02, RF-03 & Consulta y asignación de técnicos con balanceo de carga & \texttt{SVC-TEC} \\
\addlinespace
Guion, P10, P12 (sin validar) & RNF-01 & Réplicas por servicio con verificación de salud y RTO $<$ 30 min (ADR-001) & \texttt{GW-01} \\
\addlinespace
FORM-AC-01 (R-03) & RNF-02 & Mensajería asíncrona para adjuntos; respuesta \texttt{503} si proveedor falla & \texttt{SVC-INC}, cola \\
\addlinespace
FORM-AC-01 (R-01) & RF-11 / R-01 & Validación estricta de máquina de estados con \texttt{409 Conflict} & \texttt{SVC-INC} \\
\bottomrule
\end{tabularx}
\end{table}

\section{Acta de levantamiento de información}
\label{anexo:acta}

La entrevista con el proveedor de internet no se realizó antes del cierre de E1, porque el
entrevistado no estuvo disponible. Por eso no hay acta firmada ni grabación. Los requisitos citan
la pregunta del guion (\texttt{docs/levantamiento/guion-entrevista.md}) de la que se espera
obtener el dato, marcada como \emph{sin validar}, y las cifras de la línea base quedan pendientes.
La entrevista se hará antes de E2; con sus respuestas se ajustarán los requisitos y esta tabla.

\section{Verificación de referencias}

\begin{table}[H]
\centering
\caption{Comprobación de cada referencia}
\small
\begin{tabularx}{\linewidth}{@{}lYl@{}}
\toprule
Clave & Comprobación realizada & Fecha \\
\midrule
\texttt{li2022incident} & DOI (\url{https://doi.org/10.1145/3544497.3544499}) abre la obra correcta en ACM Digital Library y los datos bibliográficos coinciden con Crossref. & 22/11/2026 \\
\addlinespace
\texttt{harsh2023murphy} & DOI (\url{https://doi.org/10.1145/3603269.3604877}) abre la obra correcta en ACM SIGCOMM y metadatos verificados en Crossref. & 22/11/2026 \\
\addlinespace
\texttt{li2023identifying} & DOI (\url{https://doi.org/10.1109/ISSRE59848.2023.00028}) abre la publicación de IEEE ISSRE y metadatos verificados en Crossref. & 22/11/2026 \\
\addlinespace
\texttt{gilbert2012} & DOI (\url{https://doi.org/10.1109/MC.2011.389}) comprobado en IEEE Computer Society y verificado en Crossref. & 22/11/2026 \\
\addlinespace
\texttt{abadi2012} & DOI (\url{https://doi.org/10.1109/MC.2012.33}) comprobado en IEEE Computer Society y verificado en Crossref. & 22/11/2026 \\
\addlinespace
\texttt{basiri2016} & DOI (\url{https://doi.org/10.1109/MS.2016.60}) comprobado en IEEE Software y verificado en Crossref. & 22/11/2026 \\
\bottomrule
\end{tabularx}
\end{table}
